The mathematical study of apportionment began in earnest with a paradox.
In 1880, the chief clerk of the Census Bureau, C.\ W.\ Seaton, computed
apportionments for House sizes ranging from 275 to 350.
He discovered that when the House was sized at 299, Alabama received 8 seats;
when sized at 300, Alabama received only 7.
The addition of a seat to the House had caused a state to \emph{lose} a seat.

This result, which came to be called the \emph{Alabama Paradox},
was not a computational error.
It was a mathematical property of the Hamilton method --- the method
then in use for U.S.\ congressional apportionment.
Seaton wrote to members of Congress to report the finding,
and it triggered a political crisis: if the House size could not be
set without risking arbitrary redistributions of seats,
the entire apparatus of periodic reapportionment was in question.

The Alabama Paradox was not isolated.
Two related paradoxes were discovered in subsequent decades:
\begin{enumerate}
  \item \textbf{Population Paradox} (documented for 1900 Census):
    A state whose population grows faster than another may nonetheless
    lose a seat to the slower-growing state at the next apportionment.
  \item \textbf{New-States Paradox} (Oklahoma, 1907):
    When Oklahoma was admitted to the Union with a proportional increase
    in House size, the rounding of remainders caused Maine to gain a
    seat and New York to lose one --- an exchange that had nothing to
    do with their populations.
\end{enumerate}

All three paradoxes share a common mathematical root: Hamilton's
distribution of remainder seats by fractional-part ranking is
non-monotone.
The ranking of fractional parts changes as House size increases,
populations change, or new states are added, and these changes can
rearrange which states receive remainder seats in counterintuitive ways.

The resolution came in two stages.
First, Congress replaced Hamilton with Huntington-Hill in 1941,
eliminating the paradoxes by switching to a divisor method.
Second, Balinski and Young \citeyearpar{balinski1982fair} proved in 1982
that this tradeoff is fundamental: no apportionment method can simultaneously
satisfy the quota rule (each state within 1 of its exact quota) and be
paradox-immune.
The Alabama Paradox and the quota rule are mathematically incompatible.

\subsection{Scope and Organisation}

This paper provides:
\begin{itemize}
  \item A formal definition of each paradox and a numerical demonstration
    under Hamilton (Sections~\ref{sec:alabama}--\ref{sec:newstates}).
  \item A formal proof of divisor method immunity to all three paradoxes
    (Section~\ref{sec:immunity}).
  \item A proof sketch of the Balinski-Young impossibility theorem
    (Section~\ref{sec:balinski-young}).
  \item Policy and historical consequences (Section~\ref{sec:consequences}).
\end{itemize}
The mathematical level is accessible to readers with undergraduate
mathematics (real analysis, basic combinatorics) but sufficient for
use in expert witness testimony or appellate-level policy analysis.

\subsection{Notation}

Throughout: $n$ states, populations $p_1, \ldots, p_n > 0$,
$N = \sum_i p_i$, House size $H \geq n$,
exact quotas $q_i = p_i \times H / N$.
An apportionment is a vector $(s_1, \ldots, s_n)$ with $s_i \geq 1$
and $\sum_i s_i = H$.
For Hamilton: lower quotas $\lfloor q_i \rfloor$,
fractional parts $f_i = q_i - \lfloor q_i \rfloor$.
